% ==================================================
% PARTIE A -- Architecture et simulation LTE/EPC
% Auteur : KHLIF Iheb
% Les figures sont dans img/A/ (chemins relatifs au fichier principal).
% ==================================================

\chapter{Partie A --- Architecture et simulation LTE/EPC}

Cette partie présente l'architecture LTE/EPC telle qu'elle est modélisée dans
NS-3 (version 3.48), la simulation d'une communication de bout en bout, puis
l'influence de la distance, de la mobilité et du handover sur le débit et sur la
puissance reçue (RSRP). Les six questions du TP sont traitées dans l'ordre de
l'énoncé.

\section*{Paramètres communs aux simulations}

Sauf mention contraire, toutes les simulations de cette partie reposent sur le
même squelette, dérivé de l'exemple \texttt{lena-simple-epc.cc} (fichiers
\texttt{tp3\_q1.cc}, \texttt{tp3\_q5.cc}, \texttt{tp3\_q5\_parallel.cc} et
\texttt{tp3\_q6.cc}).

\begin{table}[H]
\centering
\caption{Paramètres communs des simulations de la partie A.}
\label{tab:A-params}
\small
\begin{tabularx}{\textwidth}{@{}lY@{}}
\toprule
\textbf{Élément} & \textbf{Valeur} \\
\midrule
Topologie & 1 eNodeB en $(0,0,0)$, 1 UE, 1 remote host relié au PGW
(liaison point-à-point de \SI{100}{Gb/s}, délai \SI{10}{ms}, MTU 1500) \\
Trafic & Deux flux UDP simultanés générés par \texttt{UdpClient} avec un
paquet toutes les \SI{1}{ms} (taille par défaut de \SI{1024}{o}) :
descendant (remote host $\rightarrow$ UE, port 1100) et montant
(UE $\rightarrow$ remote host, port 2001), démarrés à \SI{0.5}{s} \\
Charge offerte & $(1024+8+20+4)\times 8 / \SI{1}{ms} = \SI{8.448}{Mb/s}$
par sens au niveau RLC (1024 o de données, 8 o UDP, 20 o IP, environ 4 o
d'en-têtes PDCP/RLC) ; \SI{8.416}{Mb/s} au niveau IP \\
Couche radio & Paramètres par défaut du module LTE : modèle de propagation
de Friis, porteuse à \SI{2.12}{GHz}, 25 blocs de ressources (\SI{5}{MHz}),
puissance d'émission de \SI{30}{dBm} (eNodeB) et \SI{10}{dBm} (UE) \\
Traces & \texttt{lteHelper->EnableTraces()} (fichiers \texttt{DlRlcStats} et
\texttt{UlRlcStats}), \texttt{lteHelper->EnablePhyTraces()} (fichier
\texttt{DlRsrpSinrStats}) et FlowMonitor \\
\bottomrule
\end{tabularx}
\end{table}

Le choix d'une charge de \SI{8.4}{Mb/s} par sens est important pour
l'interprétation des résultats : c'est une charge \emph{fixe} (débit constant),
inférieure à la capacité maximale de la cellule. Le débit mesuré ne peut donc
jamais dépasser \SI{8.4}{Mb/s} et il ne diminue que lorsque la capacité du lien
radio devient inférieure à cette charge.

% ==================================================
\section{Question 1 --- Simulation LTE/EPC et connectivité de bout en bout}
% ==================================================

\subsection{Principe du scénario}

Le scénario construit suit la chaîne
\emph{Remote host} $\leftrightarrow$ \emph{PGW/SGW} $\leftrightarrow$
\emph{eNodeB} $\leftrightarrow$ \emph{UE}. Le cœur de réseau est instancié par
\texttt{PointToPointEpcHelper}, la partie radio par \texttt{LteHelper}. Les
étapes principales sont les suivantes :

\begin{enumerate}[nosep]
\item création des deux helpers et de la liaison entre l'EPC et le
\texttt{LteHelper} (\texttt{SetEpcHelper}) ;
\item création du remote host, installation de la pile Internet et connexion au
PGW par une liaison point-à-point ;
\item création d'un eNodeB et d'un UE, placés par un modèle de mobilité à
position constante ;
\item installation des équipements LTE, attribution d'une adresse IP à l'UE par
le PGW (\texttt{AssignUeIpv4Address}) et attachement de l'UE à l'eNodeB ;
\item configuration du routage statique : route vers le réseau \texttt{7.0.0.0/8}
(réseau des UE) sur le remote host, et passerelle par défaut de l'UE vers le
PGW ;
\item installation des applications (flux UDP, ping) et activation des traces.
\end{enumerate}

\begin{lstlisting}[language=C++]
Ptr<LteHelper> lteHelper = CreateObject<LteHelper>();
Ptr<PointToPointEpcHelper> epcHelper = CreateObject<PointToPointEpcHelper>();
lteHelper->SetEpcHelper(epcHelper);

// Routage : le remote host joint le reseau des UE via le PGW ;
// l'UE utilise le PGW comme passerelle par defaut.
remoteHostStaticRouting->AddNetworkRouteTo(Ipv4Address("7.0.0.0"),
                                           Ipv4Mask("255.0.0.0"), 1);
ueStaticRouting->SetDefaultRoute(epcHelper->GetUeDefaultGatewayAddress(), 1);

lteHelper->Attach(ueLteDevs.Get(0), enbLteDevs.Get(0));

PingHelper ping(ueIpIface.GetAddress(0));   // remote host -> UE
lteHelper->EnableTraces();
lteHelper->EnablePhyTraces();
\end{lstlisting}

\subsection{Vérification par ping}

Un ping est lancé depuis le remote host vers l'UE (adresse \texttt{7.0.0.2},
attribuée par le PGW dans le réseau \texttt{7.0.0.0/8}). Le résultat est
présenté en figure~\ref{fig:A-ping}.

\capture{img/A/q1-ping.png}{Ping du remote host vers l'UE (adresse
\texttt{7.0.0.2}) : connectivité de bout en bout vérifiée.}{fig:A-ping}

La connectivité de bout en bout est établie : les quatre requêtes reçoivent une
réponse, sans perte. Deux observations méritent d'être commentées :

\begin{itemize}[nosep]
\item \textbf{TTL = 63.} Le TTL initial (64) est décrémenté d'une unité : le
paquet a traversé un seul routeur IP, le PGW. Le tunnel GTP-U entre le PGW et
l'eNodeB est transparent pour la couche IP de bout en bout, ce qui est
précisément le rôle de l'encapsulation GTP-U.
\item \textbf{RTT d'environ \SI{29}{ms} à \SI{32}{ms}.} La liaison PGW--remote
host introduit à elle seule $2\times\SI{10}{ms}=\SI{20}{ms}$. Les \SI{9}{ms}
à \SI{12}{ms} restants correspondent à la traversée de l'interface radio (en
LTE, la transmission est organisée en sous-trames de \SI{1}{ms} ; l'ordonnancement
descendant et surtout la procédure d'allocation montante -- demande
d'ordonnancement puis attribution de ressources -- ajoutent quelques
millisecondes) ainsi qu'au traitement dans les couches PDCP/RLC/MAC. L'alternance
entre \SI{28.9}{ms} et \SI{31.9}{ms} (écart de \SI{3}{ms}) provient de
l'alignement des paquets sur la cadence des sous-trames et des opportunités
d'émission montante.
\end{itemize}

% ==================================================
\section{Question 2 --- Rôle des entités de l'architecture EPC simulée}
% ==================================================

La figure~\ref{fig:A-archi} résume l'architecture simulée. NS-3 sépare le réseau
d'accès radio (E-UTRAN, constitué des eNodeB) du cœur de réseau (EPC).

\begin{figure}[H]
\centering
\begin{tikzpicture}[
  ent/.style={draw=bleu, fill=bleu!8, rounded corners=3pt, minimum width=2.0cm,
              minimum height=1.0cm, align=center, font=\small\bfseries},
  autre/.style={ent, dashed, fill=white, font=\small},
  lien/.style={thick, bleu},
  lbl/.style={font=\scriptsize, align=center, text=black}]
  \node[ent] (ue)  at (0,0)    {UE};
  \node[ent] (enb) at (4.6,0)  {eNodeB};
  \node[ent] (gw)  at (9.2,0)  {SGW/PGW};
  \node[ent] (rh)  at (13.6,0) {Remote\\host};
  \node[ent]   (mme)  at (6.9,2.8) {MME};
  \node[autre] (enb2) at (2.0,2.8) {autre\\eNodeB};

  \draw[lien] (ue) -- (enb) node[midway, above=1pt, lbl] {Uu};
  \node[lbl, below=2pt] at (2.3,0) {PDCP/RLC/\\MAC/PHY};
  \draw[lien] (enb) -- (gw) node[midway, above=1pt, lbl] {S1-U};
  \node[lbl, below=2pt] at (6.9,0) {GTP-U/\\UDP/IP};
  \draw[lien] (gw) -- (rh) node[midway, above=1pt, lbl] {SGi (IP)};
  \node[lbl, below=2pt] at (11.4,0) {100~Gb/s\\10~ms};
  \draw[lien, dashed] (enb2) -- (enb) node[midway, above, sloped, lbl] {X2};
  \draw[lien] (enb) -- (mme) node[midway, above, sloped, lbl] {S1-MME};
  \draw[lien] (mme) -- (gw) node[midway, above, sloped, lbl] {S11};

  \draw[<->, thick, orange!80!black] (0,-1.7) -- (9.2,-1.7)
    node[midway, below, lbl] {porteuse EPS = porteuse radio (Uu) + tunnel GTP-U (S1-U)};
  \draw[<->, thick, green!50!black] (0,-2.7) -- (13.6,-2.7)
    node[midway, below, lbl] {connexion IP de bout en bout (UE \texttt{7.0.0.2} $\leftrightarrow$ remote host \texttt{1.0.0.2})};
\end{tikzpicture}
\caption{Architecture LTE/EPC simulée sous NS-3 : entités, interfaces et
piles protocolaires du plan usager.}
\label{fig:A-archi}
\end{figure}

\begin{description}[style=nextline, leftmargin=1.2em, itemsep=4pt]
\item[eNodeB (station de base).] C'est l'unique entité de l'E-UTRAN. Elle porte
la pile radio PDCP/RLC/MAC/PHY avec les UE de sa cellule (interface Uu) et gère
les ressources radio : ordonnancement (scheduler MAC), adaptation du débit
(choix du MCS selon la qualité du canal), contrôle des mesures et décision de
handover (RRC). Côté réseau, l'eNodeB relie chaque porteuse radio à un tunnel
GTP-U vers la passerelle : il est le point de jonction entre le monde radio et
le monde IP.

\item[SGW (Serving Gateway).] Passerelle de données du côté accès. Elle termine
les tunnels GTP-U venant des eNodeB et sert de point d'ancrage pour la mobilité
entre eNodeB : lors d'un handover (Q6), c'est la SGW qui redirige le tunnel
descendant vers le nouvel eNodeB sans que l'adresse IP de l'UE ne change.

\item[PGW (PDN Gateway).] Passerelle vers le réseau extérieur (ici, le remote
host). Elle attribue l'adresse IP de l'UE (réseau \texttt{7.0.0.0/8}),
encapsule dans un tunnel GTP-U les paquets descendants destinés à l'UE
(sélection de la porteuse EPS adaptée) et décapsule les paquets montants avant
de les router vers Internet. C'est elle qui apparaît comme routeur dans le
ping (TTL = 63). Dans NS-3, SGW et PGW sont regroupées dans un seul n\oe ud
(\texttt{GetPgwNode()}).

\item[MME (Mobility Management Entity).] Entité de contrôle : attachement de
l'UE, établissement des porteuses EPS et gestion de la mobilité. Dans NS-3 elle
est fortement simplifiée : la signalisation S1-MME est modélisée par des
échanges directs entre entités et non par une liaison simulée.

\item[Interface S1.] Elle relie les eNodeB à l'EPC. Sa partie plan usager,
S1-U, transporte les paquets IP de l'UE dans des tunnels GTP-U (GTP-U sur UDP
sur IP) entre l'eNodeB et la SGW/PGW ; sa partie plan de contrôle, S1-MME,
véhicule la signalisation vers la MME. Dans NS-3, S1-U est simulée par une
liaison point-à-point dédiée par eNodeB. L'interface X2, entre eNodeB voisins,
sera utilisée à la question~6 pour le handover.
\end{description}

\paragraph{Parcours d'un paquet descendant.} Le remote host émet un paquet vers
\texttt{7.0.0.2} ; le PGW le reçoit comme routeur, l'encapsule dans le tunnel
GTP-U de la porteuse EPS de l'UE ; il traverse S1-U jusqu'à l'eNodeB, qui le
décapsule puis le transmet à l'UE par les couches PDCP, RLC, MAC et PHY.
Le chemin inverse est suivi par les paquets montants. La couche IP de l'UE et
celle du remote host voient une simple connexion de bout en bout, comme le
montre la figure de la pile protocolaire de l'énoncé.

% ==================================================
\section{Question 3 --- Débit descendant et montant en fonction du temps}
% ==================================================

Les fichiers \texttt{DlRlcStats.txt} et \texttt{UlRlcStats.txt} fournissent,
par intervalle de \SI{0.25}{s}, le nombre d'octets reçus au niveau RLC. Le débit
est calculé pour chaque intervalle $[t_{d}, t_{f}]$ par
\[
  D = \frac{8 \times \text{RxBytes}}{t_{f}-t_{d}}
\]
et rapporté au milieu de l'intervalle (script \texttt{tp3\_q3.py}). La
simulation dure \SI{5}{s} avec l'UE fixe à \SI{60}{m} de l'eNodeB. Le résultat
est donné en figure~\ref{fig:A-q3}.

\capture{img/A/q3-debit-temps.png}{Débit LTE descendant (haut) et montant
(bas) en fonction du temps, d'après \texttt{DlRlcStats} et
\texttt{UlRlcStats}.}{fig:A-q3}

\paragraph{Interprétation.}
\begin{itemize}[nosep]
\item \textbf{Régime permanent.} Le débit descendant se stabilise à
\SI{8.45}{Mb/s}, valeur qui correspond exactement à la charge offerte au niveau
RLC (\SI{8.448}{Mb/s}, tableau~\ref{tab:A-params}). À \SI{60}{m} de l'eNodeB le
canal est excellent : la cellule écoule tout le trafic offert, sans perte. Le
plateau n'est donc pas la capacité de la cellule, mais la charge imposée par la
source UDP.
\item \textbf{Phase transitoire.} Le premier échantillon (environ
\SI{8.0}{Mb/s}) est plus faible : la source démarre à \SI{0.5}{s} et ses
premiers paquets n'arrivent à l'UE qu'après un délai d'une douzaine de
millisecondes (\SI{10}{ms} de liaison Internet auxquelles s'ajoute le délai
radio), ce qui retire environ 5\,\% de volume à l'intervalle
$[0.5\,;\,0.75]$~s.
\item \textbf{Variations du débit montant.} Le débit montant oscille entre
environ \SI{8.2}{Mb/s} et \SI{8.65}{Mb/s} autour de la même valeur moyenne
(\SI{8.45}{Mb/s}). Ces variations d'environ $\pm 3\,\%$ ne traduisent pas une
instabilité du lien : en montant, l'UE ne peut émettre que dans les ressources
que l'eNodeB lui attribue sous forme de blocs de transport discrets (demande
d'ordonnancement, rapport d'état du tampon, puis attribution). Les paquets
sont donc écoulés par rafales, et le nombre d'octets qui tombe dans chaque
fenêtre de \SI{0.25}{s} varie légèrement, alors que la moyenne reste égale au
débit offert. En descendant, l'ordonnanceur de l'eNodeB dispose directement du
tampon et lisse davantage la livraison.
\end{itemize}

% ==================================================
\section{Question 4 --- Influence de la distance UE--eNodeB}
% ==================================================

\subsection{Méthodologie}

L'UE est placé à position constante sur l'axe $x$ à une distance $d$ de
l'eNodeB. Le script \texttt{tp3\_q4.py} lance la simulation pour
$d \in \{100;\,300;\,500;\,1000;\,1500;\,2000;\,3000;\,4000;\,5000;\,6000;\,7000;\,8000;\,9000;\,10000\}$~m
(paramètre \texttt{--distance} de \texttt{tp3\_q1}) puis, pour chaque run :
\begin{itemize}[nosep]
\item le débit moyen DL et UL est calculé à partir de \texttt{DlRlcStats} et
\texttt{UlRlcStats} (octets reçus divisés par la durée) ;
\item le RSRP est lu dans \texttt{DlRsrpSinrStats} (colonne en watts) ; les
échantillons sont moyennés en puissance linéaire, puis convertis en dBm.
\end{itemize}

\subsection{Résultats}

\capture{img/A/q4-debit-distance.png}{Débit moyen descendant et montant en
fonction de la distance UE--eNodeB.}{fig:A-q4-debit}

\capture{img/A/q4-rsrp-distance.png}{RSRP moyen en fonction de la distance
UE--eNodeB.}{fig:A-q4-rsrp}

\begin{table}[H]
\centering
\caption{Valeurs caractéristiques relevées sur les figures~\ref{fig:A-q4-debit}
et~\ref{fig:A-q4-rsrp} (valeurs approchées) et RSRP théorique du modèle de
Friis.}
\label{tab:A-q4}
\small
\begin{tabular}{rcccc}
\toprule
Distance (m) & RSRP simulé (dBm) & RSRP Friis (dBm) & Débit DL (Mb/s) & Débit UL (Mb/s) \\
\midrule
100    & $-73{,}8$  & $-73{,}7$  & 8,42 & 8,44 \\
1000   & $-93{,}8$  & $-93{,}7$  & 8,42 & 8,44 \\
4000   & $-105{,}8$ & $-105{,}8$ & 8,42 & 8,44 \\
5000   & $-107{,}7$ & $-107{,}7$ & 8,42 & 6,43 \\
6000   & $-109{,}3$ & $-109{,}3$ & 8,42 & 4,95 \\
7000   & $-110{,}6$ & $-110{,}6$ & 7,7  & 4,95 \\
8000   & $-111{,}8$ & $-111{,}8$ & 7,2  & 4,38 \\
9000   & $-112{,}8$ & $-112{,}8$ & 7,2  & 3,49 \\
10000  & $-113{,}7$ & $-113{,}7$ & 6,18 & 3,49 \\
\bottomrule
\end{tabular}
\end{table}

\subsection{Interprétation}

\paragraph{Le RSRP suit la loi de Friis.} Le RSRP décroît de
\SI{-73.8}{dBm} à \SI{100}{m} jusqu'à \SI{-113.7}{dBm} à \SI{10}{km}, soit
\SI{20}{dB} par décade de distance : c'est la signature d'une propagation en
espace libre (exposant d'affaiblissement $n=2$). La puissance d'émission de
\SI{30}{dBm} est répartie sur $25\times12=300$ éléments de ressource, d'où une
puissance de référence par élément de $30-10\log_{10}300 \approx \SI{5.2}{dBm}$,
et
\[
  \mathrm{RSRP}(d) \approx 5{,}2 - 20\log_{10}\!\left(\frac{4\pi d f}{c}\right)
  \quad\text{avec } f=\SI{2.12}{GHz},
\]
soit $-73{,}7$~dBm à \SI{100}{m} et $-113{,}7$~dBm à \SI{10}{km}, ce qui
coïncide avec les valeurs simulées (tableau~\ref{tab:A-q4}). Cette vérification
analytique valide la chaîne de mesure (traces PHY, moyenne en puissance linéaire).
La décroissance est très rapide au voisinage de l'eNodeB (\SI{-10}{dB} entre
\SI{100}{m} et \SI{300}{m}) puis de plus en plus lente en échelle linéaire.

\paragraph{Le débit ne se dégrade qu'au-delà d'un seuil.} La
figure~\ref{fig:A-q4-debit} montre trois régimes :
\begin{enumerate}[nosep]
\item \emph{Plateau} : jusqu'à environ \SI{4}{km} dans les deux sens (et
jusqu'à \SI{6}{km} en descendant), le débit reste égal à la charge offerte
($\approx\SI{8.4}{Mb/s}$), alors que le RSRP a déjà perdu plus de \SI{30}{dB}.
Ce n'est pas contradictoire : la \emph{modulation et le codage adaptatifs}
(AMC) compensent la baisse du SINR en choisissant un MCS plus robuste (donc moins
efficace spectralement). Tant que la capacité résultante reste supérieure à
\SI{8.4}{Mb/s}, tout le trafic est écoulé.
\item \emph{Dégradation} : lorsque la capacité de la cellule devient inférieure
à la charge offerte, les tampons RLC se remplissent puis débordent ; les
paquets en excès sont perdus et le débit mesuré est égal à la capacité
disponible, qui décroît avec la distance.
\item \emph{Faible débit} : à \SI{10}{km}, il ne reste que \SI{6.2}{Mb/s} en
descendant (environ $-27\,\%$) et \SI{3.5}{Mb/s} en montant (environ
$-59\,\%$).
\end{enumerate}

\paragraph{Asymétrie montant/descendant.} Le débit montant se dégrade
nettement plus tôt (dès \SI{5}{km}) que le débit descendant (vers \SI{7}{km}).
La raison est le bilan de liaison : l'UE émet à \SI{10}{dBm} contre \SI{30}{dBm}
pour l'eNodeB, soit \SI{20}{dB} de moins. À distance égale, le SINR en réception
à l'eNodeB est donc nettement plus faible que celui reçu par l'UE : le lien
montant est le facteur limitant de la couverture, comme dans les réseaux réels.

\paragraph{Paliers.} Les courbes ne décroissent pas de façon continue mais par
paliers (par exemple \SI{4.95}{Mb/s} à \SI{6}{km} et \SI{7}{km} en montant,
\SI{7.2}{Mb/s} à \SI{8}{km} et \SI{9}{km} en descendant) : le MCS et la taille
du bloc de transport ne prennent que des valeurs discrètes.

\paragraph{Limites du modèle.} Le modèle de Friis ne comporte ni obstacles, ni
ombrage, ni évanouissements, ni interférences (une seule cellule). Il est très
optimiste : une cellule réelle en milieu urbain ne dépasse pas quelques
kilomètres, et la dégradation y serait plus rapide (exposant d'affaiblissement
de 3 à 4) et plus brutale.

% ==================================================
\section{Question 5 --- Mobilité \texttt{RandomWalk2dMobilityModel}}
% ==================================================

\subsection{Configuration}

L'UE suit désormais un modèle de marche aléatoire à vitesse constante. Il est
placé initialement à \SI{5000}{m} de l'eNodeB, distance à partir de laquelle la
qualité du lien commence à dépendre de la position (question~4), et se déplace
dans un domaine de $20\times\SI{20}{km}$ centré sur l'eNodeB. Le mode
\texttt{Distance} impose un changement de direction aléatoire tous les
\SI{2000}{m}.

\begin{lstlisting}[language=C++]
ueMobility.SetMobilityModel("ns3::RandomWalk2dMobilityModel",
    "Bounds", RectangleValue(Rectangle(-10000, 10000, -10000, 10000)),
    "Speed",  StringValue("ns3::ConstantRandomVariable[Constant=<v>]"),
    "Mode",   StringValue("Distance"),
    "Distance", DoubleValue(2000.0));
\end{lstlisting}

\subsection{Protocole expérimental}

Les directions successives de l'UE sont tirées aléatoirement : le résultat d'une
seule simulation dépend de la trajectoire tirée et n'est pas représentatif
(remarque méthodologique de l'énoncé). Chaque vitesse est donc simulée
\emph{cinq fois} avec des flux aléatoires indépendants. Le programme
\texttt{tp3\_q5\_parallel.cc} reçoit pour cela un paramètre \texttt{--run} :

\begin{lstlisting}[language=C++]
cmd.AddValue("run", "RNG run number", run);
cmd.Parse(argc, argv);
RngSeedManager::SetSeed(1);    // graine fixe
RngSeedManager::SetRun(run);   // numero de run : flux independants
\end{lstlisting}

La graine reste fixée à~1 et seul le numéro de run varie de 1 à 5 : c'est la
méthode recommandée par NS-3 pour obtenir des réplications indépendantes. Le
script \texttt{tp3\_q5\_sweep.py} compile le programme une seule fois, puis
lance les $5\ \text{vitesses}\times 5\ \text{runs}=25$ simulations de
\SI{300}{s} (vitesses 5, 10, 20, 30 et \SI{50}{m/s}), six processus au plus en
parallèle, chacune dans son propre répertoire de sortie. Les débits moyens
(octets reçus $\times 8$ divisés par la durée utile de \SI{299.5}{s}) et les
taux de perte sont obtenus avec FlowMonitor. Pour chaque vitesse, on reporte la
moyenne, l'écart-type empirique $\sigma$ (diviseur $n-1$) et la demi-largeur de
l'intervalle de confiance à 95\,\%, $t_{0{,}975;4}\,\sigma/\sqrt{5}$ avec
$t_{0{,}975;4}=2{,}776$ (loi de Student, 4 degrés de liberté).

\subsection{Évolution au cours du temps}

Les figures~\ref{fig:A-q5-dist} et~\ref{fig:A-q5-debit} donnent la distance
UE--eNodeB et les débits RLC instantanés pour une trajectoire particulière
(simulation à \SI{50}{m/s}, run~1) : les segments de \SI{2000}{m} sont
parcourus en \SI{40}{s}.

\capture{img/A/q5-distance-temps.png}{Distance UE--eNodeB au cours du temps
(RandomWalk2d, \SI{50}{m/s}, run~1 ; segments de \SI{2000}{m} parcourus en
\SI{40}{s}).}{fig:A-q5-dist}

\capture{img/A/q5-debit-temps.png}{Débit RLC descendant (haut) et montant
(bas) au cours du temps pour la trajectoire de la figure~\ref{fig:A-q5-dist}.}{fig:A-q5-debit}

On retrouve directement les résultats de la question~4 : le débit instantané ne
dépend que de la distance courante, et suit les mêmes paliers.
\begin{itemize}[nosep]
\item Au départ ($d=\SI{5}{km}$), le débit montant vaut \SI{6.45}{Mb/s}, comme
dans la figure~\ref{fig:A-q4-debit}. Il tombe à \SI{4.95}{Mb/s} vers $t=19$~s
(distance proche de \SI{6}{km}), puis à \SI{4.4}{Mb/s} lorsque l'UE atteint la
zone la plus éloignée (près de \SI{7.5}{km} à \SI{7.9}{km}).
\item Le débit descendant chute de \SI{8.45}{Mb/s} à \SI{7.7}{Mb/s} à $t\approx26$~s,
quand l'UE franchit environ \SI{6.2}{km}, et remonte à $t\approx119$~s, lorsqu'il
repasse sous cette même distance : le même seuil est donc franchi dans les deux
sens.
\item Quand l'UE se rapproche de l'eNodeB (fin de simulation, $d<\SI{4.3}{km}$),
les deux débits reviennent à la charge offerte. L'oscillation du montant en fin
de trace est celle déjà observée à la question~3.
\item À la reprise après une zone dégradée, on observe un pic bref
(\SI{8.6}{Mb/s} vers $t=119$~s) : les paquets accumulés dans les tampons sont
évacués d'un coup.
\end{itemize}

\subsection{Effet de la vitesse (moyenne sur 5 runs)}

\capture{img/A/q5-debit-vitesse.png}{Débit moyen descendant et montant en
fonction de la vitesse de l'UE (moyenne $\pm$ écart-type sur 5 runs,
FlowMonitor).}{fig:A-q5-vitesse}

\capture{img/A/q5-pertes-vitesse.png}{Taux de perte de paquets en fonction
de la vitesse de l'UE (moyenne $\pm$ écart-type sur 5 runs).}{fig:A-q5-pertes}

\begin{table}[H]
\centering
\caption{Débit moyen sur \SI{300}{s} en fonction de la vitesse de l'UE :
moyenne, écart-type $\sigma$ et demi-largeur de l'IC à 95\,\% sur 5 runs
(Mb/s).}
\label{tab:A-q5}
\small
\begin{tabular}{rcccccc}
\toprule
 & \multicolumn{3}{c}{Descendant} & \multicolumn{3}{c}{Montant} \\
\cmidrule(lr){2-4}\cmidrule(lr){5-7}
Vitesse (m/s) & Moyenne & $\sigma$ & IC\,95\,\% & Moyenne & $\sigma$ & IC\,95\,\% \\
\midrule
5  & 8,38 & 0,05 & $\pm 0{,}07$ & 6,16 & 0,28 & $\pm 0{,}34$ \\
10 & 8,18 & 0,16 & $\pm 0{,}20$ & 5,70 & 0,29 & $\pm 0{,}36$ \\
20 & 8,07 & 0,23 & $\pm 0{,}29$ & 5,53 & 0,49 & $\pm 0{,}61$ \\
30 & 7,98 & 0,42 & $\pm 0{,}52$ & 5,55 & 0,80 & $\pm 0{,}99$ \\
50 & 7,79 & 0,61 & $\pm 0{,}75$ & 5,50 & 1,22 & $\pm 1{,}52$ \\
\bottomrule
\end{tabular}
\end{table}

\begin{table}[H]
\centering
\caption{Taux de perte de paquets en fonction de la vitesse de l'UE : moyenne,
écart-type $\sigma$ et demi-largeur de l'IC à 95\,\% sur 5 runs (\%).}
\label{tab:A-q5-pertes}
\small
\begin{tabular}{rcccccc}
\toprule
 & \multicolumn{3}{c}{Descendant} & \multicolumn{3}{c}{Montant} \\
\cmidrule(lr){2-4}\cmidrule(lr){5-7}
Vitesse (m/s) & Moyenne & $\sigma$ & IC\,95\,\% & Moyenne & $\sigma$ & IC\,95\,\% \\
\midrule
5  & 0,3 & 0,5 & $\pm 0{,}6$ & 25,7 & 3,0  & $\pm 3{,}8$ \\
10 & 2,6 & 1,8 & $\pm 2{,}3$ & 30,9 & 3,4  & $\pm 4{,}2$ \\
20 & 4,0 & 2,6 & $\pm 3{,}2$ & 33,2 & 5,4  & $\pm 6{,}7$ \\
30 & 4,8 & 4,6 & $\pm 5{,}7$ & 32,9 & 8,8  & $\pm 10{,}9$ \\
50 & 7,1 & 7,0 & $\pm 8{,}6$ & 33,4 & 13,7 & $\pm 17{,}0$ \\
\bottomrule
\end{tabular}
\end{table}

\begin{table}[H]
\centering
\caption{Résultats individuels des 25 simulations : débit descendant / débit
montant (Mb/s).}
\label{tab:A-q5-runs}
\footnotesize
\begin{tabular}{rccccc}
\toprule
Vitesse (m/s) & Run 1 & Run 2 & Run 3 & Run 4 & Run 5 \\
\midrule
5  & 8,33 / 5,88 & 8,31 / 5,85 & 8,42 / 6,42 & 8,42 / 6,42 & 8,42 / 6,23 \\
10 & 8,01 / 5,40 & 8,00 / 5,39 & 8,31 / 6,00 & 8,27 / 5,94 & 8,31 / 5,79 \\
20 & 7,86 / 5,01 & 8,20 / 5,84 & 8,09 / 5,52 & 7,82 / 5,11 & 8,36 / 6,19 \\
30 & 8,04 / 5,39 & 8,18 / 5,82 & 8,03 / 5,41 & 7,28 / 4,45 & 8,38 / 6,66 \\
50 & 8,19 / 6,08 & 7,94 / 5,39 & 7,54 / 4,81 & 6,86 / 4,02 & 8,39 / 7,22 \\
\bottomrule
\end{tabular}
\end{table}

\paragraph{Stabilité de la connexion.} Dans aucune des 25 simulations la
connexion n'est rompue : l'UE reste attaché à l'unique eNodeB et les débits
moyens restent supérieurs à \SI{6.8}{Mb/s} en descendant et à \SI{4.0}{Mb/s}
en montant (tableau~\ref{tab:A-q5-runs}), ce qui exclut toute coupure
prolongée. La dégradation est progressive et réversible, conformément au
modèle de propagation sans évanouissements. Elle se manifeste par les pertes de
paquets, en moyenne de \SI{0.3}{\%} à \SI{7.1}{\%} en descendant et de
\SI{26}{\%} à \SI{33}{\%} en montant selon la vitesse ; les valeurs extrêmes
d'un run atteignent \SI{17.9}{\%} en descendant et \SI{50.3}{\%} en montant
(\SI{50}{m/s}, run~4), à comparer à l'absence totale de perte à courte distance.

\paragraph{Tendance moyenne.} Contrairement à ce que donnerait une exécution
unique, la moyenne sur 5 runs décroît de façon régulière avec la vitesse
(figures~\ref{fig:A-q5-vitesse} et~\ref{fig:A-q5-pertes}) : le débit descendant
passe de \SI{8.38}{Mb/s} à \SI{7.79}{Mb/s} (environ $-7\,\%$) quand la vitesse
passe de 5 à \SI{50}{m/s}, et les pertes descendantes de \SI{0.3}{\%} à
\SI{7.1}{\%}. En montant, le débit passe de \SI{6.16}{Mb/s} à \SI{5.50}{Mb/s}
(environ $-11\,\%$), mais la baisse est concentrée entre 5 et \SI{20}{m/s} : au-delà,
les moyennes sont indiscernables (\SIrange{5.50}{5.55}{Mb/s}, pertes autour de
\SI{33}{\%}). À titre de comparaison, le run~1 seul
(tableau~\ref{tab:A-q5-runs}, première colonne) donne
\num{8.33}, \num{8.01}, \num{7.86}, \num{8.04} puis \SI{8.19}{Mb/s} en
descendant : une courbe en « V » non monotone, qui aurait conduit à une
conclusion erronée si une seule trajectoire avait été simulée.

\paragraph{Pourquoi la moyenne diminue.} NS-3 ne modélise pas, par défaut,
l'effet Doppler ni les évanouissements rapides : la qualité du lien ne dépend
que de la \emph{distance instantanée} (figure~\ref{fig:A-q5-debit}). La
vitesse n'agit donc que sur l'étendue de l'espace explorée en \SI{300}{s}. À
\SI{5}{m/s}, l'UE ne parcourt que \SI{1.5}{km}, donc moins d'un segment : il
reste dans un anneau de \SI{5}{km} à environ \SI{6.5}{km} où la qualité varie
peu. À \SI{50}{m/s}, il parcourt \SI{15}{km} (sept segments) et explore tout le
domaine. L'UE part en bordure de la zone où le débit est égal à la charge
offerte : une trajectoire qui le rapproche de l'eNodeB ne peut pas relever le
débit au-delà de \SI{8.4}{Mb/s} (plafond fixé par la source), alors qu'une
trajectoire qui l'en éloigne le fait chuter par paliers (question~4). Plus
l'UE explore, plus l'asymétrie « gain plafonné, perte non bornée » tire la
moyenne vers le bas. Cette explication est cohérente avec la saturation observée
en montant, où l'UE franchit déjà le seuil de dégradation à \SI{5}{km}, mais
reste une interprétation du modèle et non un résultat mesuré directement.

\paragraph{Variabilité entre graines.} L'écart-type augmente fortement avec
la vitesse : il est multiplié par environ 11 en descendant ($0{,}05$ à
$\SI{0.61}{Mb/s}$) et par 4 en montant ($0{,}28$ à $\SI{1.22}{Mb/s}$). À faible
vitesse, la trajectoire est quasiment déterminée par la première direction
tirée ; à grande vitesse, le résultat dépend de toute la ligne brisée, qui
diffère fortement d'un run à l'autre. Le tableau~\ref{tab:A-q5-runs} le
montre : le run~5 reste excellent à toutes les vitesses (descendant entre
\num{8.31} et \SI{8.42}{Mb/s} et montant en hausse jusqu'à \SI{7.22}{Mb/s} à
\SI{50}{m/s}, signe d'une trajectoire qui revient vers l'eNodeB), alors que le
run~4 se dégrade régulièrement (descendant de \num{8.42} à \SI{6.86}{Mb/s}, montant de
\num{6.42} à \SI{4.02}{Mb/s}). À \SI{5}{m/s}, les runs~3 et~4 donnent des
valeurs identiques : sur un seul segment de \SI{1.5}{km}, les deux directions
tirées maintiennent l'UE dans la même zone de qualité (même palier de MCS).

\paragraph{Variabilité comparée à l'effet de la vitesse.} L'effet moyen de la
vitesse reste du même ordre que la dispersion entre graines. Entre 5 et
\SI{50}{m/s}, la baisse moyenne du débit descendant est de \SI{0.59}{Mb/s},
alors que l'écart-type à \SI{50}{m/s} est de \SI{0.61}{Mb/s} et que l'écart
entre le meilleur et le pire run à cette vitesse atteint \SI{1.53}{Mb/s}
(soit environ 2,6 fois l'effet moyen). Les intervalles de confiance à 95\,\%
des vitesses 5 et \SI{50}{m/s} ($8{,}38 \pm 0{,}07$~Mb/s et
$7{,}79 \pm 0{,}75$~Mb/s) se recouvrent. Des tests de Welch (cinq valeurs par
vitesse) donnent, en descendant, $p\approx0{,}05$ entre 5 et \SI{10}{m/s},
$p\approx0{,}04$ entre 5 et \SI{20}{m/s} et $p\approx0{,}10$ entre 5 et
\SI{50}{m/s} ; une analyse de variance sur les cinq vitesses n'est pas
significative ($p\approx0{,}14$ en descendant, $p\approx0{,}58$ en montant).
En montant, seule la différence entre \SI{5}{m/s} et les vitesses de 10 et
\SI{20}{m/s} est à la limite de la significativité ($p\approx0{,}03$ à $0{,}05$) ;
entre 20 et \SI{50}{m/s}, aucune différence n'est décelable ($p\approx0{,}96$).
Avec seulement cinq runs par configuration et sans correction pour les
comparaisons multiples, ces tests ont une faible puissance : la décroissance
du débit avec la vitesse est cohérente et plausible, mais elle n'est
établie que de façon modérée. Pour la confirmer, il faudrait davantage de
runs, d'autant que la dispersion entre graines croît avec la vitesse. En
conditions réelles, une vitesse plus élevée dégraderait en plus le lien (effet
Doppler, évanouissements rapides, handovers plus fréquents), ce que ce scénario
à une seule cellule ne reproduit pas.

% ==================================================
\section{Question 6 --- Deux eNodeB, interface X2 et handover}
% ==================================================

\subsection{Scénario}

Un second eNodeB est ajouté. Les deux stations sont placées en
$(0,0,0)$ et $(500,0,0)$ ; l'UE démarre en $(-100,0,0)$ et se déplace en ligne
droite à \SI{10}{m/s} le long de l'axe $x$ (\texttt{ConstantVelocityMobilityModel})
pendant \SI{100}{s}. Sa trajectoire traverse donc successivement les deux
cellules : il passe à la verticale de l'eNodeB~1 à $t=10$~s, du point
équidistant ($x=250$~m) à $t=35$~s, puis de l'eNodeB~2 à $t=60$~s. L'interface X2
est activée avec \texttt{AddX2Interface()} et le handover automatique est
déclenché par l'algorithme \texttt{A2A4RsrqHandoverAlgorithm}, paramétré comme
dans l'énoncé :

\begin{lstlisting}[language=C++, basicstyle=\ttfamily\footnotesize]
lteHelper->SetHandoverAlgorithmType("ns3::A2A4RsrqHandoverAlgorithm");
lteHelper->SetHandoverAlgorithmAttribute("ServingCellThreshold",
                                         UintegerValue(30));
lteHelper->SetHandoverAlgorithmAttribute("NeighbourCellOffset",
                                         UintegerValue(1));
...
lteHelper->AddX2Interface(enbNodes);
lteHelper->Attach(ueLteDevs.Get(0), enbLteDevs.Get(0)); // initial: eNodeB 1

Config::ConnectWithoutContext(
    "/NodeList/*/DeviceList/*/LteEnbRrc/HandoverStart",
    MakeCallback(&NotifyHandoverStart));
Config::ConnectWithoutContext(
    "/NodeList/*/DeviceList/*/LteUeRrc/HandoverEndOk",
    MakeCallback(&NotifyHandoverEndOk));
\end{lstlisting}

Les callbacks \texttt{HandoverStart} et \texttt{HandoverEndOk} (traces RRC)
affichent l'instant du handover ; la distance de l'UE aux deux eNodeB et le
débit (par fenêtres de \SI{100}{ms}, via FlowMonitor) sont enregistrés en
parallèle.

\subsection{Résultats}

Le handover se produit à $t=\SI{35.84}{s}$ : l'UE quitte l'eNodeB~1 pour
l'eNodeB~2. Il n'y en a qu'un seul sur toute la trajectoire. Même après avoir
dépassé l'eNodeB~2 ($t>60$~s), l'UE reste attaché à celui-ci, car il en est
toujours plus proche que de l'eNodeB~1.

\capture{img/A/q6-distances.png}{Distances de l'UE aux deux eNodeB. La
droite verticale marque l'instant du handover ($t=\SI{35.84}{s}$).}{fig:A-q6-dist}

\capture{img/A/q6-rsrp.png}{RSRP mesuré par l'UE sur sa cellule serveuse
(eNodeB~1 puis eNodeB~2) ; le handover a lieu à $t=\SI{35.84}{s}$.}{fig:A-q6-rsrp}

\capture{img/A/q6-debit-handover.png}{Débits descendant et montant autour du
handover (fenêtres de \SI{100}{ms}).}{fig:A-q6-debit}

\begin{table}[H]
\centering
\caption{Synthèse du handover.}
\label{tab:A-q6}
\small
\begin{tabular}{@{}ll@{}}
\toprule
Instant du handover & \SI{35.84}{s} (position de l'UE : $x=\SI{258.4}{m}$) \\
Distance à l'eNodeB 1 / eNodeB 2 & \SI{258.4}{m} / \SI{241.6}{m} \\
RSRP de la cellule source (juste avant) & $\approx -82{,}0$~dBm \\
RSRP de la cellule cible (juste après) & $\approx -81{,}4$~dBm \\
Débit descendant (nominal $\to$ minimum) & $\approx 8{,}4 \to 8{,}2$~Mb/s ($\approx -3\,\%$) \\
Débit montant (nominal $\to$ minimum) & $\approx 8{,}4 \to 4{,}6$~Mb/s ($\approx -45\,\%$) \\
Durée de la perturbation & une seule fenêtre de \SI{100}{ms} touchée (moins de \SI{200}{ms}) \\
\bottomrule
\end{tabular}
\end{table}

\subsection{Interprétation}

\paragraph{Instant du handover.} Le point équidistant des deux stations
($x=\SI{250}{m}$) est atteint à $t=35$~s, alors que le handover n'est déclenché
qu'à $t=\SI{35.84}{s}$, soit \SI{0.84}{s} (environ \SI{8}{m}) plus tard. Ce
retard est voulu : l'algorithme A2-A4 sur le RSRQ n'agit pas dès que la cellule
voisine est meilleure. Il faut que la qualité (RSRQ) de la cellule serveuse
passe sous le seuil \texttt{ServingCellThreshold} (événement A2), puis que celle
d'une voisine dépasse la cellule serveuse d'au moins \texttt{NeighbourCellOffset}
(événement A4), et ces mesures sont lissées et rapportées périodiquement. Cette
hystérésis évite l'effet « ping-pong » (bascules répétées) près de la frontière
des cellules. À l'instant du handover, la cellule cible est reçue à
$\approx-81{,}5$~dBm contre $\approx-82{,}0$~dBm pour la cellule source (valeurs
de Friis aux distances de \SI{241.6}{m} et \SI{258.4}{m}) : la cible est à peine
meilleure, mais le handover est déjà déclenché.

\paragraph{Impact sur le RSRP.} Le RSRP de la cellule serveuse (figure~\ref{fig:A-q6-rsrp})
est continu à la bascule : la cellule~1 est suivie à $\approx-82$~dBm
puis la cellule~2 à $\approx-81{,}4$~dBm, soit un saut de moins de
\SI{1}{dB}. Cela confirme que le handover est déclenché dans la zone où les deux
cellules sont équivalentes ; l'UE ne subit donc aucune dégradation radio du fait
de la bascule. Les deux pics à $+5$~dBm ($t=10$~s et $t=60$~s) sont des
artefacts du scénario : la trajectoire passe exactement à la verticale des
antennes ($d=0$), où le modèle de Friis sature sans perte de propagation. Ils
n'ont pas de sens physique (pas de hauteur d'antenne dans le modèle) et ne
doivent pas être interprétés. Après le handover, le RSRP de la cellule~2 augmente
jusqu'au pic de $t=60$~s puis décroît ($\approx-85{,}7$~dBm à $t=100$~s,
$d=\SI{390}{m}$).

\paragraph{Impact sur le débit.} La figure~\ref{fig:A-q6-debit} montre que la
bascule ne provoque qu'une perturbation très brève : une seule fenêtre de
\SI{100}{ms} est touchée, avec un creux d'environ $-3\,\%$ en descendant et de
près de $-45\,\%$ en montant (\SI{4.6}{Mb/s}), immédiatement suivi d'un retour
au nominal avec un léger dépassement (\SI{8.85}{Mb/s}) correspondant à
l'évacuation des paquets retardés. L'asymétrie entre les deux sens s'explique
par la procédure de handover : les paquets descendants en attente dans
l'eNodeB source sont retransmis à l'eNodeB cible via X2, tandis que, côté
montant, l'UE ne peut recommencer à émettre qu'après son accès aléatoire (RACH)
sur la cellule cible. La précision de cette mesure est limitée par la fenêtre
d'observation de \SI{100}{ms} : la vraie durée d'interruption est de l'ordre de
quelques dizaines de millisecondes à une centaine de millisecondes.

\paragraph{Remarque pour la suite du TP.} Pour la comparaison avec la bascule
WiFi/LTE de la partie~B (question~4), ce handover intra-LTE fournit la valeur de
référence suivante : interruption inférieure à \SI{200}{ms}, creux montant
d'environ 45\,\% (donc légèrement au-dessus du seuil de 50\,\% de la
définition de l'énoncé) et quasi nul en descendant.

% ==================================================
\section{Synthèse de la partie A}
% ==================================================

\begin{itemize}[nosep]
\item \textbf{Architecture.} La connectivité de bout en bout est vérifiée (ping,
TTL = 63) ; l'EPC simulé se compose de la SGW/PGW (regroupées), d'une MME
simplifiée et des interfaces S1-U (GTP-U) et X2.
\item \textbf{Débit en régime nominal.} \SI{8.45}{Mb/s} dans les deux sens,
égal à la charge offerte : la cellule n'est pas saturée.
\item \textbf{Distance.} Le RSRP suit la loi de Friis (\SI{20}{dB} par décade).
L'AMC maintient le débit jusqu'à environ \SI{4}{km} (\SI{6}{km} en descendant) ;
au-delà, le débit chute par paliers, plus tôt en montant à cause de la
puissance d'émission de l'UE.
\item \textbf{Mobilité.} Sans évanouissements ni effet Doppler, le débit ne
dépend que de la distance instantanée ; la vitesse n'agit que sur l'étendue de
l'espace parcouru. Sur 5 runs par vitesse, le débit moyen décroît légèrement
avec la vitesse (descendant : \num{8.38} à \SI{7.79}{Mb/s} ; montant :
\num{6.16} à \SI{5.50}{Mb/s}), mais la dispersion entre graines, qui croît avec
la vitesse, est du même ordre que cet effet : la tendance n'est que
modérément établie statistiquement.
\item \textbf{Handover.} Un seul handover automatique (A2-A4 sur RSRQ), déclenché à
\SI{35.84}{s} avec un léger retard par rapport au point équidistant, sans effet
notable sur le RSRP et avec une interruption inférieure à \SI{200}{ms} (surtout
visible en montant).
\end{itemize}
